%%%%%%%%%%%%%%%%%%%%%%%%%%%%%%%%%%%%%%%%%%%%%%%%%%%%%%%%%%%%%%%%%%%%%%%%%%%%%%
%%  notations.tex
%%  The 3/2 Phase Transition in Semi-Random 2-SAT
%%
%%  Symbols are organised by SCOPE, as variables are in a program:
%%
%%    Part I    general machinery (delimiters, operators, utilities)
%%    Part II   probability and expectation
%%    Part III  GLOBAL symbols: defined in Section 1, valid everywhere
%%    Part IV   MODULE symbols of Section 2 (graph-theoretic analysis)
%%    Part V    MODULE symbols of Section 3 (the main identity)
%%    Part VI   MODULE symbols of Section 4 and Appendices A, B
%%
%%  Rules.
%%   1. A global symbol is never reused for anything else.  Every symbol that
%%      appears in a statement of Section 1 is global.
%%   2. A module symbol is valid only in its own section (and appendices).  It
%%      may not reuse the letter of a global symbol.  Two different modules may
%%      reuse a letter, since their scopes do not overlap.
%%   3. A symbol that lives inside one proof (a local symbol) is written by
%%      hand in that proof and is not defined here.  It may not reuse the
%%      letter of a global symbol or of a module symbol of the same module.
%%   4. A subscript or a prime may specialise a symbol but not change its
%%      meaning: R_0, R_1 and R' are radii, p_0 is a retention probability.
%%   5. A symbol is defined only if it is used at least twice, or appears in
%%      a statement.  Everything else is written out in words.
%%
%%  checkscope.py checks rules 1--4 mechanically.  To rename a global or
%%  module symbol, edit the single line that defines it and recompile.
%%%%%%%%%%%%%%%%%%%%%%%%%%%%%%%%%%%%%%%%%%%%%%%%%%%%%%%%%%%%%%%%%%%%%%%%%%%%%%


%%%%%%%%%%%%%%%%%%%%%%%%%%%%%%%%%%%%%%%%%%%%%%%%%%%%%%%%%%%%%%%%%%%%%%%%%%%%%%
%% Part I.  General machinery
%%%%%%%%%%%%%%%%%%%%%%%%%%%%%%%%%%%%%%%%%%%%%%%%%%%%%%%%%%%%%%%%%%%%%%%%%%%%%%

% Delimiters that grow with their content in display style but stay small
% inline.  Everything below is built from these, so the paper never needs a
% hand-written \left...\right.
\newcommand{\AutoAdjust}[3]{%
  \mathchoice{\left#1#2\right#3}{#1#2#3}{#1#2#3}{#1#2#3}}
\newcommand{\InParentheses}[1]{\AutoAdjust{(}{#1}{)}}
\newcommand{\InBrackets}[1]{\AutoAdjust{[}{#1}{]}}
\newcommand{\InBraces}[1]{\AutoAdjust{\{}{#1}{\}}}

\newcommand{\bk}[1]{\InParentheses{#1}}
\newcommand{\Bk}[1]{\InBrackets{#1}}
\newcommand{\BK}[1]{\InBraces{#1}}
\newcommand{\abs}[1]{\left|{#1}\right|}
\newcommand{\norm}[1]{\left\lVert#1\right\rVert}
\newcommand{\ceil}[1]{\left\lceil{#1}\right\rceil}
\newcommand{\floor}[1]{\left\lfloor{#1}\right\rfloor}
\newcommand{\angbk}[1]{\left\langle{#1}\right\rangle}

% Sets and operators
\newcommand{\Z}{\mathbb{Z}}
\newcommand{\R}{\mathbb{R}}
\newcommand{\N}{\mathbb{N}}
\newcommand{\eps}{\varepsilon}
\newcommand{\dd}{\mathrm{d}}
\DeclareMathOperator*{\argmin}{arg\,min}
\DeclareMathOperator*{\argmax}{arg\,max}

\newcommand{\eqdef}{\overset{\mathrm{def}}{=\mathrel{\mkern-3mu}=}}
\newcommand{\defeq}{\eqdef}

% Utilities
\newcommand{\bnote}[1]{\textcolor{blue}{[#1]}}
\newcommand{\numberthis}{\addtocounter{equation}{1}\tag{\theequation}}
\newenvironment{proofof}[1]{\begin{proof}[Proof of #1]}{\end{proof}}


%%%%%%%%%%%%%%%%%%%%%%%%%%%%%%%%%%%%%%%%%%%%%%%%%%%%%%%%%%%%%%%%%%%%%%%%%%%%%%
%% Part II.  Probability and expectation
%%
%%  \E and \Prob are the bare symbols, for use as operators: \E\cyc{\Akept}.
%%  \Ex{...} and \Prx{...} carry their own delimiters and grow with content.
%%  Both take an optional subscript: \Ex[\Phi]{...}, \Prx[\mathrm{nb}]{...}.
%%
%%  To switch the whole paper between Pr(...) and Pr[...], or between an
%%  upright Pr and \mathbb P, edit only the four lines marked [STYLE].
%%%%%%%%%%%%%%%%%%%%%%%%%%%%%%%%%%%%%%%%%%%%%%%%%%%%%%%%%%%%%%%%%%%%%%%%%%%%%%

\newcommand{\E}{\mathbb{E}}                     % [STYLE] expectation symbol
\DeclareMathOperator{\Prob}{Pr}                 % [STYLE] probability symbol
% The optional argument is a subscript; \ifblank keeps an empty one from
% typesetting an empty subscript box.  The \! cancels the thin space TeX puts
% between the symbol and the delimiter, which after an operator such as Pr is
% wide enough to read as a gap.
\newcommand{\Ex}[2][]{\E\ifblank{#1}{}{_{#1}}\!\InBrackets{#2}}        % [STYLE] E[ ... ]
\newcommand{\Prx}[2][]{\Prob\ifblank{#1}{}{_{#1}}\!\InParentheses{#2}} % [STYLE] Pr( ... )
\newcommand{\Var}[2][]{\operatorname{\mathbf{Var}}\ifblank{#1}{}{_{#1}}\InBrackets{#2}}

% Indicator.  \1 is kept as a short alias for the bare symbol, so that the
% existing subscript form \1_{\{...\}} continues to work.
\newcommand{\1}{\mathbbm{1}}
\newcommand{\indic}[1]{\1\InBraces{#1}}

% Named events and distributions
\newcommand{\Bin}{\operatorname{Bin}}
\newcommand{\event}[1]{\mathcal{#1}}


%%%%%%%%%%%%%%%%%%%%%%%%%%%%%%%%%%%%%%%%%%%%%%%%%%%%%%%%%%%%%%%%%%%%%%%%%%%%%%
%% Part III.  GLOBAL symbols
%%
%%  Defined in Section 1 and listed in its notation table
%%  (Section~\ref{subsec:notation}).  Valid from the first page to the last.
%%%%%%%%%%%%%%%%%%%%%%%%%%%%%%%%%%%%%%%%%%%%%%%%%%%%%%%%%%%%%%%%%%%%%%%%%%%%%%

%% --- The input and its random formula -------------------------------------
\newcommand{\Gin}{G}                        % the input finite multigraph
\newcommand{\Vset}{V}                       % its vertex set
\newcommand{\Eset}{E}                       % its edge set
\newcommand{\nvert}{n}                      % |V|
\newcommand{\medge}{m}                      % |E|
\newcommand{\bvar}[1]{x_{#1}}               % Boolean variable at a vertex
\newcommand{\rform}[1]{\Phi(#1)}            % the random 2-CNF on a graph
\newcommand{\fml}{\varphi}                  % a fixed (non-random) 2-CNF formula
\DeclareMathOperator{\Sol}{Sol}             % solution set of a formula
\newcommand{\SAT}{\text{ SAT}}              % the event "... is satisfiable"
\newcommand{\indsat}[1][G]{\1_{_{\text{SAT}}}(\Phi(#1))}

%% --- Graph functionals ----------------------------------------------------
\newcommand{\ncomp}[1]{c(#1)}               % number of connected components
\newcommand{\cyc}[1]{\nu(#1)}               % cycle rank |E|-|V|+c
\newcommand{\exc}[1]{\xi(#1)}               % excess |E| - (3/2)|V|
\newcommand{\dg}[2]{\deg_{#1}(#2)}          % degree of a vertex
\newcommand{\dist}{\operatorname{dist}}     % graph distance

%% --- Balls ----------------------------------------------------------------
\newcommand{\ball}[2]{B_{#1}(#2)}           % radius-#1 ball around #2
\newcommand{\ballsz}[1]{M_{#1}}             % bound on the size of a radius-#1 ball

%% --- Solution graph -------------------------------------------------------
%  Sol(phi) is viewed as a graph (assignments at Hamming distance 1 are
%  joined), so c(.) has one meaning only: components of a graph.
\newcommand{\solcomp}[1]{\ncomp{\Sol(#1)}}  % number of solution clusters
\newcommand{\freeset}{S}                    % the set of free variables of a subcube
\newcommand{\fixasg}{\tau}                  % the values of the fixed variables
\newcommand{\cplx}[1]{X(#1)}                % cubical complex of subcubes inside Sol
\newcommand{\euler}[1]{\chi(#1)}            % Euler characteristic

%% --- The random edge-deletion experiment ----------------------------------
\newcommand{\Tdef}{T}                       % number of deferred edges
\newcommand{\Kdef}{K}                       % forest components holding a deferred edge
\newcommand{\Ffor}{F}                       % the forest built in the first stage

%% --- Percolation and long paths -------------------------------------------
\newcommand{\perc}[2]{#1_{#2}}              % H_p: keep each edge independently w.p. p
\newcommand{\keepp}{p}                      % retention probability
\newcommand{\plen}{\ell}                    % length of the paths we track
\newcommand{\plenz}{\ell_0}                 % the length finally selected
\newcommand{\Lpath}[3]{\mathcal{L}_{#1}(#2,#3)} % P(random root starts a length-#1 path)

%% --- Parameters -----------------------------------------------------------
\newcommand{\excpar}{\varepsilon}           % the excess parameter, xi(G) >= eps n
\newcommand{\Hyp}{\textup{(H)}}             % the standing hypothesis
% \loctree carries its own math delimiters and is used in TEXT mode, so that
% the words can be broken across lines.
\newcommand{\loctree}[2]{\((#1,#2)\)-locally tree-like}
\newcommand{\degb}{\Delta}                  % maximum-degree bound
\newcommand{\excd}{\delta}                  % excess density, xi(G) >= delta n
%  R is THE letter for a radius.  R and rho are the parameters of
%  Definition 1.1; R_0, R_1 are particular radii, R' is a radius inside a proof.
\newcommand{\rad}{R}                        % radius in the local-tree condition
\newcommand{\radz}{R_0}                     % threshold for R (Theorem 1.12)
\newcommand{\radp}{R_1}                     % radius of the packed balls (Lemmas 1.7, 1.8)
\newcommand{\nontree}{\rho}                 % tolerated fraction of non-tree balls
\newcommand{\nontreez}{\rho_0}              % threshold used for the bound on T
\newcommand{\nontreeo}{\rho_1}              % threshold of Theorem 1.12
\newcommand{\rate}{\gamma}                   % E[(3/4)^T] <= 2 e^{-gamma n}
\newcommand{\pack}{\eta}                    % fraction of vertices with a rich ball


%%%%%%%%%%%%%%%%%%%%%%%%%%%%%%%%%%%%%%%%%%%%%%%%%%%%%%%%%%%%%%%%%%%%%%%%%%%%%%
%% Part IV.  MODULE symbols of Section 2
%%
%%  None.  Every symbol of Section 2 is either global or local to one proof.
%%%%%%%%%%%%%%%%%%%%%%%%%%%%%%%%%%%%%%%%%%%%%%%%%%%%%%%%%%%%%%%%%%%%%%%%%%%%%%


%%%%%%%%%%%%%%%%%%%%%%%%%%%%%%%%%%%%%%%%%%%%%%%%%%%%%%%%%%%%%%%%%%%%%%%%%%%%%%
%% Part V.  MODULE symbols of Section 3
%%%%%%%%%%%%%%%%%%%%%%%%%%%%%%%%%%%%%%%%%%%%%%%%%%%%%%%%%%%%%%%%%%%%%%%%%%%%%%

\newcommand{\face}[2]{Q(#1,#2)}             % subcube (face) of the Boolean cube
\newcommand{\edgefac}[2]{w_{#1}(#2)}        % edge factor w_S(e): 3/4, 1/2 or 0


%%%%%%%%%%%%%%%%%%%%%%%%%%%%%%%%%%%%%%%%%%%%%%%%%%%%%%%%%%%%%%%%%%%%%%%%%%%%%%
%% Part VI.  MODULE symbols of Section 4 and Appendices A, B
%%
%%  The constants below are produced in exactly this order, and none of them
%%  depends on R or rho, which are chosen last.
%%%%%%%%%%%%%%%%%%%%%%%%%%%%%%%%%%%%%%%%%%%%%%%%%%%%%%%%%%%%%%%%%%%%%%%%%%%%%%

\newcommand{\Akept}{{\Ffor^{+}}}            % F plus the kept deferred edges (each edge w.p. 1/2)
\newcommand{\keeppz}{p_0}                   % the retention probability of Lemma 4.2
\newcommand{\plow}{\zeta}                   % lower bound on the long-path functional
\newcommand{\thr}{\kappa}                   % the error term of Lemma 4.3
\newcommand{\cycd}{\omega}                  % E nu(G_{1/2}) >= omega n

%% --- Appendix A: flows on paths -------------------------------------------
% \flowt carries its own math delimiters and is used in TEXT mode.
\newcommand{\flowt}[2]{\((#1,#2)\)-flow}
\newcommand{\brnum}{\lambda}                % branching rate demanded of good paths
\newcommand{\gfrac}{\psi}                   % fraction of good starting edges
\newcommand{\drift}{\mu}                    % mean drift of the log-branching walk
\newcommand{\ubd}{U}                        % upper bound on that walk's increments
\newcommand{\Pset}[1]{\mathcal{P}_{\plen}(#1)}  % candidate paths from a starting edge
\newcommand{\wt}[1]{w_{#1}}                 % the path weights
\newcommand{\pref}{\pi}                     % a path prefix

%% --- Appendix B: extinction in forests ------------------------------------
\newcommand{\sep}{\mathrm{Sep}}             % the separator set in a forest
\newcommand{\extp}[1]{h(#1)}                % forest extinction bound
